\documentclass[11pt]{amsart}
\usepackage{amssymb}
\usepackage[margin=1.15in]{geometry}
\usepackage[hidelinks]{hyperref}
\hypersetup{pdftitle={Descent of finite freeness}}

\newtheorem{theorem}{Theorem}[section]
\newtheorem{lemma}[theorem]{Lemma}
\newtheorem{proposition}[theorem]{Proposition}
\newtheorem{corollary}[theorem]{Corollary}
\theoremstyle{definition}
\newtheorem{definition}[theorem]{Definition}
\newcommand{\Spec}{\operatorname{Spec}}
\newcommand{\Hom}{\operatorname{Hom}}
\newcommand{\Bas}{\mathcal B}
\newcommand{\Frac}{\mathcal F}
\newcommand{\tr}[1]{\lVert #1\rVert}
\newcommand{\loc}[1]{\lVert #1\rVert_T}
\newcommand{\Span}{\operatorname{Span}}
\newcommand{\GL}{\operatorname{GL}}

\title{Descent of finite freeness}
\date{October 23, 2026}

\begin{document}
\begin{abstract}
Let $M$ be an $R$-module whose underlying type is an \'etale sheaf.
We prove that if $M$ is \'etale-locally free of a fixed finite rank, then
it merely has a basis of that rank. We also prove descent of a basis
together with a linear retraction from a finite free ambient module.
\end{abstract}
\maketitle

\section{Setting and statement}

We work in synthetic algebraic geometry over $R$, with locality,
synthetic quasi-coherence, and Zariski-local choice as in \cite{SAG}.
In particular, $R$ is a commutative local ring, $1\ne0$, and for every
finitely presented commutative $R$-algebra $A$, evaluation is an isomorphism
\[
 A\longrightarrow (\Spec A\longrightarrow R).
\]
Here $\Spec A$ is the type of $R$-algebra homomorphisms $A\to R$.

All existence assertions are understood constructively: $\exists z$
asserts that a witness merely exists. The notation $\tr Z$ likewise
asserts that $Z$ is inhabited, without choosing an element of $Z$.
An isomorphism of modules always means an $R$-linear isomorphism.
The underlying type of a module is a set, so its equalities are propositions.

For a monic polynomial $f\in R[t]$ of degree $d>0$, put
\[
 B_f=R[t]/(f),\qquad X_f=\Spec B_f=\{x\in R:f(x)=0\},
 \qquad G_f=\tr{X_f}.
\]
The topology $T$ is generated by the propositions $G_f$. A type $Y$ is
a \emph{sheaf} if the constant map
\[
 Y\longrightarrow (G_f\longrightarrow Y)
\]
is an equivalence for every $f$. For a proposition $P$, this says precisely
that $(G_f\to P)\to P$ for every $f$.

We assume propositional sheafification: $\loc Z$ is a proposition and a
sheaf, receives a map from $Z$, and any map from $Z$ to a sheaf
proposition factors through $\loc Z$. Thus a sheaf proposition may be
proved under a hypothesis which holds $T$-locally.

Write $\Bas_r(M)$ for the type of ordered bases $(v_1,\ldots,v_r)$ of
$M$, or equivalently of linear isomorphisms $R^r\to M$.

\begin{theorem}\label{thm:descent}
Let $r\geq0$ be fixed and let $M$ be an $R$-module whose underlying type
is a sheaf. Then
\[
 \loc{\Bas_r(M)}\longrightarrow\tr{\Bas_r(M)}.
\]
\end{theorem}

We shall use the following results about schemes. They follow from the
SAG axioms and do not involve descent of modules.
\begin{enumerate}
\item Zariski-local choice on an affine $X=\Spec A$ gives finitely many
  $g_i\in A$ generating the unit ideal, together with actual choices on
  each $D(g_i)$, from a family of merely inhabited types on $X$.
  The same conclusion, with a finite cover by open affine charts, holds
  on a scheme.
\item If $P_n(x)$ is an increasing sequence of propositions on a scheme
  $X$, then
  \[
   (\forall x\in X\;\exists n\;P_n(x))
   \ \Longrightarrow\ \exists N\;\forall x\in X\;P_N(x).
  \]
  This is the uniform bound theorem for schemes. Intersections of open
  charts are schemes, so it applies to these intersections as well.
\item An open predicate on an affine scheme is a finite union of standard
  opens. A finite family $D(g_i)$ covers $\Spec A$ exactly when the $g_i$
  generate $1$ in $A$. Open propositions are stable under double negation.
\end{enumerate}
The fraction argument below is the argument for sections in
\cite[Section~7.1]{SAG}. We include it to specify exactly what is used.

Sheaves are closed under products, dependent products, dependent sums,
and equality types. Consequently, if $Y$ and $Z$ are sheaves and
$u:Y\to Z$ is fixed, then ``$u$ is an equivalence'' is a sheaf
proposition. Indeed, express it by a left inverse and a right inverse
with their equations; the function types and equality types in this
description are sheaves.

The ring $R$ is a sheaf. Here is the coefficient argument. The algebra
$B_f$ is free on $1,t,\ldots,t^{d-1}$, and evaluation identifies it with
$R^{X_f}$. If $h:G_f\to R$, let $b\in B_f$ represent
$x\mapsto h(|x|)$. For each $x\in X_f$, the element $b$ is equal to
$h(|x|)\cdot1$: evaluate at any $y\in X_f$ and use that $G_f$ is a
proposition. The constant coefficient of $b$ is therefore a value of
$h$. Uniqueness follows from injectivity of $R\to B_f$ and of evaluation.

\begin{lemma}\label{lem:constant}
Let $Y$ be a set which is a sheaf, and put $X=X_f$.
If $u:X\to Y$ satisfies $u(x)=u(y)$ for all $x,y\in X$, there is a
unique $v\in Y$ with $v=u(x)$ for every $x\in X$.
Also, an equality between elements of a sheaf follows if it holds for
every $x\in X$.
\end{lemma}
\begin{proof}
Given $q:\tr X$, the common value of $u$ exists and is unique, so it can
be chosen uniquely. This gives a map $\tr X\to Y$, which has a unique
constant value because $Y$ is a sheaf. For the second assertion use that
equality types of a sheaf are sheaves.
\end{proof}

\section{Fractions of module-valued functions}

For $a\in R$, the notation $D(a)$ means the proposition that $a$ is a
unit. For a function $g:X\to R$, it means the subtype of $X$ where
$g(x)$ is a unit.

\begin{definition}\label{def:fractions}
An $R$-module $L$ satisfies the \emph{fraction condition}, denoted
$\Frac(L)$, if the following hold for every $a\in R$.
\begin{enumerate}
\item For every $s:D(a)\to L$, there exist $n\geq0$ and $u\in L$ such
  that $s=u/a^n$ on $D(a)$.
\item If $u\in L$ is zero on $D(a)$, there exists $n\geq0$ such that
  $a^nu=0$.
\end{enumerate}
Both clauses assert mere existence. In the terminology of
\cite[Definition~7.1.5]{SAG}, this is weak quasi-coherence.
\end{definition}

The condition $\Frac(L)$ is a proposition. It is preserved by linear
isomorphisms and finite products. For products, choose finitely many
exponents and raise them to a common exponent.
The module $R$ satisfies it: use evaluation for $R[1/a]$, together with
the elementary fraction and kernel descriptions of this localization.
Hence every finite free module satisfies it.

\begin{lemma}[Fractions of sections]\label{lem:sections}
Let $X$ be a scheme and let $L_x$ be an $R$-module for each $x\in X$,
with $\Frac(L_x)$ for every $x$. For $g:X\to R$, the following hold.
\begin{enumerate}
\item Every section $s$ of $L_x$ over $D(g)$ can be written
  \[
   s(x)=u(x)/g(x)^n
  \]
  for some $n\geq0$ and some section $u$ on all of $X$.
\item If a section $u$ on $X$ vanishes on $D(g)$, then
  $g^nu=0$ on $X$ for some $n\geq0$.
\end{enumerate}
In particular, the module $\prod_{x\in X}L_x$ satisfies the fraction
condition over $R$.
\end{lemma}
\begin{proof}
At each $x$, the first fraction clause gives a numerator and an exponent
for $s(x)$. The assertion that an exponent at most $n$ suffices is
increasing in $n$. The uniform bound theorem therefore gives one
exponent $n$ which suffices everywhere. Zariski-local choice then gives
a finite open cover $U_i$ of $X$ and numerators $u_i$ on $U_i$, with
$u_i=g^ns$ on $U_i\cap D(g)$.

On $U_i\cap U_j$, the difference $u_i-u_j$ vanishes on $D(g)$. The
second fraction clause, followed by the uniform bound theorem on this
intersection, gives an exponent which kills that difference everywhere
on the intersection. Choose a common exponent $k$ for the finitely many
pairs. The sections $g^ku_i$ now agree on overlaps, and hence give a
section $u$ on $X$. On $D(g)$ it satisfies $u=g^{n+k}s$.

For the second assertion apply the second fraction clause at every
$x$, and then apply the uniform bound theorem to $g(x)^nu(x)=0$.
Taking $g$ constant and exchanging the section arguments proves the
last assertion.
\end{proof}

\begin{lemma}[Kernels]\label{lem:kernel}
If $P$ and $Q$ satisfy the fraction condition and $u:P\to Q$ is linear,
then $\ker u$ satisfies the fraction condition.
\end{lemma}
\begin{proof}
Let $s:D(a)\to\ker u$. In $P$, write $s=p/a^n$. The element $u(p)$
vanishes on $D(a)$, so $a^ku(p)=0$ for some $k$. Thus $a^kp\in\ker u$
and $s=(a^kp)/a^{n+k}$. The second fraction clause is inherited from $P$.
\end{proof}

\begin{proposition}\label{prop:fractions-descend}
If the underlying type of $M$ is a sheaf, then the proposition
$\Frac(M)$ is a sheaf.
\end{proposition}
\begin{proof}
Fix $f$ and put $X=X_f$. Suppose $G_f\to\Frac(M)$. In particular, for
each $x\in X$ we have $\Frac(M)$. Lemma~\ref{lem:sections} gives the
fraction condition for
\[
 P=M^X,\qquad Q=P^X=M^{X\times X}.
\]
Here we apply the lemma first to obtain the condition for $P$, and then
again to obtain it for $Q$. We prove the two clauses for $M$.

Let $a\in R$ and $s:D(a)\to M$. Apply the first clause for $P$ to the
function which is constant in $x$. There exist $n\geq0$ and $b:X\to M$
such that
\[
 s=b(x)/a^n\qquad\text{on }D(a),\quad x\in X.
\]
Hence $b(x)-b(y)$ vanishes on $D(a)$. Apply the second clause for $Q$
to this function of $(x,y)$. For some $k\geq0$,
\[
 a^k(b(x)-b(y))=0\qquad(x,y\in X).
\]
The function $x\mapsto a^kb(x)$ is constant, so
Lemma~\ref{lem:constant} supplies $m\in M$ with $m=a^kb(x)$ for all
$x$. On $D(a)$, the equality
\[
 s=m/a^{n+k}
\]
holds for every $x\in X$. Since equality in $M$ is a sheaf, it holds
without the hypothesis $x\in X$.

For the second clause, suppose $m\in M$ vanishes on $D(a)$. Apply the
second clause for $P$ to the constant function with value $m$. It gives
$a^nm=0$ for every $x\in X$, with one exponent $n$. Descend this equality
by Lemma~\ref{lem:constant}.
\end{proof}

\begin{corollary}\label{cor:local-free-fractions}
If $M$ is a sheaf and $\loc{\Bas_r(M)}$ holds, then $\Frac(M)$ holds.
Thus Lemma~\ref{lem:sections} applies to the constant family $M$ on
every scheme.
\end{corollary}
\begin{proof}
The proposition $\Frac(M)$ is a sheaf, so we may prove it locally,
where $M\cong R^r$.
\end{proof}

\section{Choosing solutions of linear equations}

The following affine gluing result is the additive statement of
\cite[Theorem~7.3.6]{SAG}. We give its fraction proof.

\begin{lemma}\label{lem:additive-gluing}
Let $X=\Spec A$ be affine and let $K_x$ be modules satisfying the fraction
condition. Suppose $E_x$ is a merely inhabited set on which addition by
$K_x$ is defined, with
\[
 e+0=e,\qquad(e+k)+l=e+(k+l),
\]
and, for every $e,e'\in E_x$, exactly one $k\in K_x$ satisfies
$e+k=e'$. Then there merely exists a section $x\mapsto e(x)\in E_x$.
\end{lemma}
\begin{proof}
By Zariski-local choice, choose $g_1,\ldots,g_p\in A$ generating $1$,
and sections $e_i$ on $D(g_i)$. On an overlap write
\[
 \delta_{ij}=e_i-e_j\in K_x.
\]
These differences satisfy
$\delta_{ij}+\delta_{jk}=\delta_{ik}$ on triple overlaps.
By Lemma~\ref{lem:sections}, and finite choice of a common exponent,
\[
 \delta_{ij}=\frac{b_{ij}}{(g_ig_j)^n}
 \quad\text{on }D(g_ig_j)
\]
for sections $b_{ij}$ of $K_x$ on $X$ and one $n\geq0$.
The sections
\[
 g_i^n b_{jk}-g_j^n b_{ik}+g_k^n b_{ij}
\]
vanish on $D(g_ig_jg_k)$. Apply the second clause of
Lemma~\ref{lem:sections}, and choose a common exponent $q$ for all
triples. Put
\[
 F_i=g_i^{n+q},\qquad \mu_{ij}=(g_ig_j)^q b_{ij}.
\]
We then have, on all of $X$,
\begin{equation}\label{eq:cocycle}
 F_i\mu_{jk}-F_j\mu_{ik}+F_k\mu_{ij}=0.
\end{equation}
The elements $F_i$ still generate $1$. Indeed, powers of finitely many
elements generating $1$ again generate $1$, by expanding a sufficiently
high power of an identity $\sum_i a_ig_i=1$. Choose $\rho_i\in A$ with
$\sum_i\rho_iF_i=1$, and set
\[
 V_i=\sum_k\rho_k\mu_{ik}.
\]
Equation~\eqref{eq:cocycle} gives
\[
 F_jV_i-F_iV_j=\sum_k\rho_kF_k\mu_{ij}=\mu_{ij}.
\]
On $D(g_i)$ let $t_i=V_i/F_i$. On an overlap,
\[
 t_i-t_j=\frac{\mu_{ij}}{F_iF_j}=\delta_{ij}.
\]
Consequently the sections $e_i-t_i$ agree on overlaps and give a section
of $E_x$ on $X$.
\end{proof}

\begin{lemma}\label{lem:linear-choice}
Let $u:P\to Q$ be linear, with $P$ and $Q$ satisfying the fraction
condition. Let $X$ be affine and $q:X\to Q$. If
\[
 \forall x\in X\;\exists p\in P\quad u(p)=q(x),
\]
then there merely exists $p:X\to P$ with $u(p(x))=q(x)$ for all $x$.
\end{lemma}
\begin{proof}
Put $K=\ker u$, which satisfies the fraction condition by
Lemma~\ref{lem:kernel}. For each $x$, take
\[
 E_x=\{p\in P:u(p)=q(x)\}.
\]
Addition by $K$ preserves $E_x$, and the difference of two elements of
$E_x$ is the unique element of $K$ taking one to the other. Apply
Lemma~\ref{lem:additive-gluing}.
\end{proof}

\section{Coefficients and membership in a span}

In this section $M$ is a sheaf and we fix
\begin{equation}\label{eq:local-free}
 \loc{\Bas_r(M)}.
\end{equation}
The assumptions remain in force until the end of the proof of
Theorem~\ref{thm:descent}.

\begin{lemma}\label{lem:coefficients}
For $f$ monic of degree $d>0$ and $X=X_f$, the map
\[
 E_M:M^d\longrightarrow M^X,\qquad
 (m_0,\ldots,m_{d-1})\longmapsto
 \left(x\longmapsto\sum_{k=0}^{d-1}x^k m_k\right)
\]
is a linear isomorphism.
\end{lemma}
\begin{proof}
The source and target are sheaves, so the assertion that this fixed map
is an equivalence is a sheaf proposition. We may therefore use
\eqref{eq:local-free} and take $M=R^r$. In each coordinate the assertion
is evaluation
\[
 R^d\cong R[t]/(f)\cong R^X.
\]
Linearity is immediate from the formula.
\end{proof}

For $h:X\to M$, define $c_M(h)\in M$ to be the coefficient of $x^0$ in
the unique expansion in Lemma~\ref{lem:coefficients}. This map is linear
and has the two properties
\begin{equation}\label{eq:constant-coefficient}
 c_M(x\mapsto m)=m,\qquad
 c_M(x\mapsto a(x)m)=c_R(a)m.
\end{equation}
For the first, the expansion is $(m,0,\ldots,0)$. For the second, if
$a(x)=\sum_k x^ka_k$, then $a(x)m=\sum_k x^k(a_km)$.

\begin{proposition}\label{prop:span}
For fixed $w_1,\ldots,w_s,m\in M$, the proposition
\[
 m\in\Span_R(w_1,\ldots,w_s)
 \quad\Longleftrightarrow\quad
 \exists a_1,\ldots,a_s\in R\quad m=\sum_i a_iw_i
\]
is a sheaf.
\end{proposition}
\begin{proof}
Fix $f$ and suppose this proposition holds under $G_f$. It then holds
for each $x\in X_f$. Apply Lemma~\ref{lem:linear-choice} on $X_f$ to
\[
 R^s\longrightarrow M,\qquad(a_i)\longmapsto\sum_i a_iw_i,
\]
and the constant function with value $m$. The source satisfies the
fraction condition, and so does $M$ by
Corollary~\ref{cor:local-free-fractions}.
We obtain functions $a_i:X_f\to R$ such that
$m=\sum_i a_i(x)w_i$ for every $x$. Apply $c_M$ to this equality and
use \eqref{eq:constant-coefficient}:
\[
 m=\sum_i c_R(a_i)w_i.
\]
This proves the required implication $(G_f\to P)\to P$.
\end{proof}

\section{Selecting a basis}

An open proposition is one which can be expressed as a finite disjunction
of assertions that specified elements of $R$ are units.
Finite unions and finite intersections of open propositions are open.

\begin{lemma}\label{lem:matrices}
For a matrix over $R$, surjectivity is an open proposition. For a square
matrix, surjectivity is equivalent to invertibility. If
$w_1,\ldots,w_s$ generate $R^r$, then some $r$ of them form a basis.
\end{lemma}
\begin{proof}
For a matrix with columns $w_i=(w_i',a_i)\in R^{r-1}\times R$, a solution
of $\sum_i\lambda_iw_i=(0,1)$ gives $\sum_i\lambda_i a_i=1$.
By locality, some $a_j=u$ is a unit. Put
\[
 z_i=u w_i'-a_iw_j'\in R^{r-1}.
\]
If the original matrix is surjective, the $z_i$ generate $R^{r-1}$:
lift $(u^{-1}y,0)$ to obtain a combination of the $z_i$ equal to $y$.
Conversely, if the $z_i$ generate, then for a prescribed vector
$(v',v_r)$ first solve
\[
 \sum_i c_i z_i=u v'-v_rw_j'.
\]
Writing $w=\sum_i c_iw_i$, the vector
\[
 w+u^{-1}(v_r-w_r)w_j
\]
equals $(v',v_r)$. Thus surjectivity is equivalent to the finite
disjunction, over $j$, of
\[
 a_j\text{ is a unit}\quad\text{and}\quad(z_i)_i
 \text{ is surjective onto }R^{r-1}.
\]
Induction on $r$ proves openness.

For a square matrix, put a unit entry of the last row in the last
column and eliminate the other entries of that row. Surjectivity then
reduces to surjectivity of the square upper-left block. By induction
that block is invertible; the resulting block triangular matrix is
invertible. This also includes the empty matrix.

Finally, apply induction to choose $r-1$ of the $z_i$ forming a basis.
Take the corresponding $w_i$, together with $w_j$. The preceding
formula solves every equation in these $r$ columns. Hence their square
matrix is invertible, and they form a basis. For $r=0$ use the empty basis.
\end{proof}

We recall why open propositions and their families descend for this
topology. First, $\neg\neg X_f$: otherwise evaluation would identify
$B_f$ with the one-element ring, contrary to its free basis containing
$1$ and $1\ne0$ in $R$. If $G_f\to P$ and $P$ is open, then
$\neg\neg P$, and stability of open propositions gives $P$.
Thus every open proposition is a sheaf.

For an open predicate $U(x)$ on $X_f=\Spec B_f$, the proposition
$\forall x\;U(x)$ is open. To see this, write $U$ as a finite union of
$D(g_i)$. This union covers precisely when the $g_i$ generate $1$ in
$B_f$. Equivalently,
\[
 B_f^p\longrightarrow B_f,\qquad(b_i)\longmapsto\sum_i g_ib_i
\]
is surjective. In the power basis this is surjectivity of a finite
matrix over $R$, which is open by Lemma~\ref{lem:matrices}.

It follows that the type of open propositions, identified under logical
equivalence, is a sheaf. Indeed, for a family of open propositions $P_q$
indexed by $q\in G_f$, take the open proposition
\[
 P=\forall x\in X_f\;P_{|x|}.
\]
For each $q\in G_f$, we have $P\leftrightarrow P_q$: one direction uses
that $G_f$ is a proposition, and the other eliminates $q:\tr{X_f}$ into
the proposition $P_q$. Uniqueness follows because implications between
open propositions are sheaves.

\begin{lemma}\label{lem:open-local}
If $Q$ is a sheaf proposition, then ``$Q$ is open'' is a sheaf proposition.
\end{lemma}
\begin{proof}
Suppose $G_f\to(Q\text{ is open})$. Regard this as a family of open
propositions, each with underlying proposition $Q$. The preceding result
gives an open proposition $P$ with $G_f\to(P\leftrightarrow Q)$.
Both implications $P\to Q$ and $Q\to P$ are sheaves, so they hold
without $G_f$. Hence $P\leftrightarrow Q$, and $Q$ is open.
\end{proof}

\begin{lemma}\label{lem:basis-among-generators}
Under \eqref{eq:local-free}, any finite generating family in $M$ merely
contains a basis of length $r$.
\end{lemma}
\begin{proof}
For a fixed tuple $v=(v_1,\ldots,v_r)$ in $M$, the assertion that $v$ is
a basis is a sheaf proposition: it says that the fixed linear map
$R^r\to M$ is an equivalence. By Lemma~\ref{lem:open-local}, its openness
may be proved locally. There $M\cong R^r$, so it is invertibility of a
square matrix, an open condition by Lemma~\ref{lem:matrices}.

Now let $w_1,\ldots,w_s$ generate $M$. The proposition that some tuple
$(w_{i_1},\ldots,w_{i_r})$ is a basis is a finite disjunction of open
propositions, and hence is open and a sheaf. We may prove it locally
using \eqref{eq:local-free}. For $M\cong R^r$ it follows from
Lemma~\ref{lem:matrices}.
\end{proof}

\section{Finitely many generators and the descent theorem}

\begin{proposition}\label{prop:generators}
Continue to assume \eqref{eq:local-free}. Fix a monic $f$ of positive
degree. If $X_f\to\tr{\Bas_r(M)}$, then $M$ merely has finitely many
generators.
\end{proposition}
\begin{proof}
Put $X=X_f$ and $d=\deg f$. Zariski-local choice gives finitely many
$g_i\in B_f$ generating $1$, and, on each $D(g_i)$, a basis
\[
 v_{i1}(x),\ldots,v_{ir}(x)\quad\text{of }M.
\]
By Corollary~\ref{cor:local-free-fractions} and
Lemma~\ref{lem:sections}, there are functions $w_{ij}:X\to M$ and
exponents $n_i$ such that
\[
 w_{ij}(x)=g_i(x)^{n_i}v_{ij}(x)\quad\text{on }D(g_i).
\]
There are finitely many columns and charts, so these choices and a common
exponent within each chart are legitimate finite choices.
By Lemma~\ref{lem:coefficients}, write
\[
 w_{ij}(x)=\sum_{k=0}^{d-1}x^k m_{ijk},\qquad m_{ijk}\in M.
\]
We claim that the finitely many elements $m_{ijk}$ generate $M$.

Fix $m\in M$ and let $P(m)$ assert membership in their span. This is a
sheaf proposition by Proposition~\ref{prop:span}. For each $x\in X$,
some $g_i(x)$ is a unit, since the $g_i$ generate $1$ and $R$ is local.
Every $w_{ij}(x)$ is in the proposed span by its displayed expansion.
Multiplying by $g_i(x)^{-n_i}$ shows that all the $v_{ij}(x)$ are in the
span. They form a basis, so $P(m)$ holds. Thus $X\to P(m)$, and hence
$G_f\to P(m)$. Since $P(m)$ is a sheaf, $P(m)$ holds globally.
\end{proof}

\begin{proof}[Proof of Theorem~\ref{thm:descent}]
Fix \eqref{eq:local-free} and put $Q=\tr{\Bas_r(M)}$.
We show that $Q$ is a sheaf. Given $G_f\to Q$, we have $X_f\to Q$.
Proposition~\ref{prop:generators} supplies finitely many generators of
$M$, and Lemma~\ref{lem:basis-among-generators} supplies a basis of
length $r$. Hence $(G_f\to Q)\to Q$ for every $f$, as required.

We can now eliminate \eqref{eq:local-free} into the sheaf proposition
$Q$, using the map which sends a basis to its mere existence.
This yields $\tr{\Bas_r(M)}$.
\end{proof}

The order of this proof matters. The assumption
$\loc{\Bas_r(M)}$ first proves the fraction condition, the coefficient
isomorphism, descent of span membership, and openness of a fixed basis
condition. Each of those targets is already known to be a sheaf
proposition. Finitely many generators are then constructed under the
separate hypothesis $X_f\to\tr{\Bas_r(M)}$. Only after this construction
do we prove that mere existence of a basis is itself a sheaf.

\section{Bases with a retraction}

\begin{lemma}\label{lem:left-inverse}
Let $F:R^r\to R^m$ be fixed and linear. The proposition
\[
 \exists H:R^m\to R^r\text{ linear}\qquad HF=1
\]
is a sheaf.
\end{lemma}
\begin{proof}
Write $F_{ki}$ for the matrix of $F$. For each $j=1,\ldots,r$, the
$j$th row $(h_{j1},\ldots,h_{jm})$ of a left inverse must satisfy
\[
 \sum_{k=1}^m h_{jk}F_{ki}=\delta_{ji}\qquad(1\leq i\leq r).
\]
Solvability is membership of the vector $(\delta_{ji})_i$ in the span of
the fixed vectors $(F_{ki})_i$, for $1\leq k\leq m$. It is a sheaf
proposition by Proposition~\ref{prop:span} applied to the already free
module $R^r$. The conjunction of the $r$ row conditions is a sheaf.
Finite choice assembles solutions of the rows into $H$.
\end{proof}

\begin{theorem}\label{thm:split}
Let $E$ be a merely finite free $R$-module and let $S\subseteq E$ be a
submodule such that membership in $S$ is a sheaf proposition. If
$T$-locally there are an isomorphism $b_0:R^r\to S$ and a linear map
$h_0:E\to R^r$ with $h_0\iota b_0=1$, where $\iota:S\hookrightarrow E$,
then such a pair merely exists.
\end{theorem}
\begin{proof}
We may choose a basis of $E$, since the conclusion asserts mere
existence. The underlying type of $E$ is a sheaf, and so is
$S=\{e\in E:e\in S\}$. Forgetting $h_0$ gives local freeness of $S$;
Theorem~\ref{thm:descent} gives a basis $b:R^r\to S$.

Fix this basis. Each local pair $(b_0,h_0)$ gives a left inverse of
$\iota b$ by the formula
\[
 h=b^{-1}b_0h_0:E\longrightarrow R^r.
\]
Indeed, $h_0\iota=b_0^{-1}$ because $b_0$ is an isomorphism, and hence
$h\iota b=1$. After identifying $E$ with $R^m$, existence of this left
inverse is a sheaf proposition by Lemma~\ref{lem:left-inverse}.
Its local existence therefore gives mere existence of $h$ globally.
The pair $(b,h)$ has the required properties.
\end{proof}

\begin{corollary}
If an $R$-algebra $A$ has a sheaf as its underlying type and is
$T$-locally isomorphic to $M_N(R)$, then its underlying module merely
has a basis of rank $N^2$. For a right ideal with sheaf membership,
local existence of a rank-$r$ basis with a linear retraction from $A$
therefore implies mere existence of such a basis and retraction.
\end{corollary}
\begin{proof}
The underlying module of $M_N(R)$ is free of rank $N^2$. Apply
Theorem~\ref{thm:descent} to the underlying module of $A$, and then apply
Theorem~\ref{thm:split} to the right ideal as a submodule.
\end{proof}

\begin{thebibliography}{9}
\bibitem{SAG} F.~Cherubini, T.~Coquand, M.~Hutzler,
\emph{A foundation for synthetic algebraic geometry},
arXiv:2307.00073.
\end{thebibliography}
\end{document}
